\documentclass[../../../main.tex]{subfiles}

\begin{document}

\section{Quick refresher on the current in electrodynamics}
\label{sec:current-electrodynamics}

\subsection{Maxwell's equations}

Given Maxwell's equations,
\begin{align}
\vec\nabla\cdot\vec E(\vec r,t)&=\frac{1}{\varepsilon_0}\,\rho(\vec r,t) \label{eq:maxwell_a}\\
\vec\nabla\cdot\vec B(\vec r,t)&=0 \label{eq:maxwell_b}\\
\vec\nabla\times\vec E(\vec r,t)&=-\frac{\partial}{\partial t}\vec B(\vec r,t) \label{eq:maxwell_c}\\
\vec\nabla\times\vec B(\vec r,t)&=\frac{1}{c^2}\frac{\partial}{\partial t}\vec E(\vec r,t)
   +\frac{1}{\varepsilon_0 c^2}\,\vec j(\vec r,t) \label{eq:maxwell_d}
\end{align}

\subsection{Continuity equation}

From $\varepsilon_0\,\partial_t\,(\cref{eq:maxwell_a})$, using $\partial_x\partial_t=\partial_t\partial_x$,
\begin{equation}
\frac{\partial}{\partial t}\rho(\vec r,t)
=\varepsilon_0\,\vec\nabla\cdot\bigl(\partial_t\vec E(\vec r,t)\bigr),
\end{equation}
and from $\varepsilon_0 c^2\,\vec\nabla\cdot(\cref{eq:maxwell_d})$,
\begin{align}
\vec\nabla\cdot\vec j(\vec r,t)
&=\vec\nabla\cdot\bigl(\vec\nabla\times\vec B(\vec r,t)\bigr)
-\varepsilon_0\,\vec\nabla\cdot\bigl(\partial_t\vec E(\vec r,t)\bigr)\notag\\
&= -\varepsilon_0\,\vec\nabla\cdot\bigl(\partial_t\vec E(\vec r,t)\bigr) \qquad (\operatorname{div}(\operatorname{rot}\vec F)=0).
\end{align}
Adding the two gives the continuity equation
\begin{equation}
\frac{\partial}{\partial t}\rho(\vec r,t)+\vec\nabla\cdot\vec j(\vec r,t)=0.\label{eq:continuity}
\end{equation}

\subsection{Point-particle solutions}

Solutions of the continuity equation:
\begin{align}
\rho(\vec r,t)&=\sum_\alpha q_\alpha\,\delta\bigl(\vec r-\vec r_\alpha(t)\bigr) \label{eq:continuity_solution_rho}\\
\vec j(\vec r,t)&=\sum_\alpha q_\alpha\,\vec v_\alpha(t)\,\delta\bigl(\vec r-\vec r_\alpha(t)\bigr)
\label{eq:continuity_solution_j}
\end{align}

\begin{proof}
Note that
\begin{align}
\partial_t\,\delta\bigl(\vec r-\vec r_\alpha(t)\bigr)
&=\partial_t\Bigl(\delta\bigl(r_x-r_x^\alpha(t)\bigr)\delta\bigl(r_y-r_y^\alpha(t)\bigr)
\delta\bigl(r_z-r_z^\alpha(t)\bigr)\Bigr)\notag\\
&= -\sum_{i=x,y,z}\dot r_i^\alpha(t)\,\partial_{r_i}\delta\bigl(r_i-r_i^\alpha(t)\bigr)
=-\dot{\vec r}_\alpha(t)\cdot\vec\nabla\,\delta\bigl(\vec r-\vec r_\alpha(t)\bigr),
\end{align}
where we used the chain rule in the second line and
\begin{align}
\vec\nabla\cdot\Bigl(\vec v_\alpha(t)\,\delta\bigl(\vec r-\vec r_\alpha(t)\bigr)\Bigr)
&=\sum_{i=x,y,z}\partial_{r_i}\Bigl(v_i^\alpha(t)\,\delta\bigl(\vec r-\vec r_\alpha(t)\bigr)\Bigr)\notag\\
&=\sum_{i=x,y,z}v_i^\alpha(t)\,\partial_{r_i}\delta\bigl(\vec r-\vec r_\alpha(t)\bigr)
=\vec v_\alpha(t)\cdot\vec\nabla\,\delta\bigl(\vec r-\vec r_\alpha(t)\bigr).
\end{align}
Therefore
\begin{equation}
\frac{\partial}{\partial t}\rho(\vec r,t)
=-\sum_\alpha q_\alpha\,\dot{\vec r}_\alpha(t)\cdot
\Bigl(\vec\nabla\,\delta\bigl(\vec r-\vec r_\alpha(t)\bigr)\Bigr)
=-\vec\nabla\cdot\vec j(\vec r,t),
\end{equation}
which is \eqref{eq:continuity}.
\end{proof}

Note that these solutions are not unique (one equation for four components). 
\paragraph{Intermediate result: the current can be defined by ``charge times velocity''.}

\end{document}
